\Article{356}
\textsc{Theorema.} \emph{Si $P$ sit functio ipsius $x$ prima $m$ dimensionum et $X$ functio ipsorum $x,x^p,x^{pp},x^{p^3}\dotstwo x^{p^{m-1}}$, in quam omnes hae quantitates aequaliter ingrediantur, i. e. quae eadem maneat, quomodocunque eae inter se permutentur, functio $X$ per $P$ divisa dabit residuum, quod erit numerus.}

\emph{Dem.} Sit residuum
\[
Ax^{m-1}+Bx^{m-2}+\dotstwo+N\equiv\xi
\]
omnes coëfficientes $A,B,C\ldots$ usque ad $N$ exclusive erunt $\equiv0$. Hoc ita demonstratur. Quum $X-\xi$ per $P$ dividatur, etiam $X^p-\xi^p$ per $P$ dividi 
\SourceLeaf{238}
poterit. Sed facile perspicitur, $X^p$ esse id, quod fit $X$, si pro $x$ ponatur $x^p$, pro $x^p$, $x^{pp}$ etc.\ .\ et pro $x^{p^{m-1}}$, $x^{p^m}$ seu quod idem est $x$. Hinc patet, esse $X^p\equiv X\;(\mathrm{mod.}\,P)$; quare, quum $X^p\equiv\xi^p$ et $X\equiv\xi(\mathrm{mod.}\,P)$, erit etiam $\xi^p\equiv\xi(\mathrm{mod.}\,P$ seu
\[
\xi^p-\xi\equiv0(\mathrm{mod.}\,P)
\]
At $\xi^p-\xi$ secundum modulum $p$ congruum est producto ex $\xi,\xi+1,\xi+2,\dotstwo$ usque ad $\xi+p-1$, qui factores omnes ad $P$ primi erunt, nisi $\xi$ sit simpliciter numerus. Quare etiam $\xi^p-\xi$ alio modo per $P$ divisibilis non erit. Q. E. D.

Huiusmodi functiones sunt summa omnium, summa quadratorum, cuborum etc., summa productorum e binis, ternis etc. Quis vero sit ille numerus, per \textsection\ sq. determinabimus:

\Article{357}
\textsc{Theorema.} \emph{Sit functio prima \textsection\ praec.}
\[
P\equiv x^m-Ax^{m-1}+Bx^{m-2}-Cx^{m-3}+\text{etc.}
\]
\emph{erit residuum, si summa quantitatum $x,x^p$ etc. $x^{p^{m-1}}$ per $P$ dividatur, $\equiv A$, si summa productorum e binis, $\equiv B$, si summa productorum e ternis, $\equiv C$ etc.}

\emph{Dem.} Sint functiones illae $X,Y,Z$ etc. earumque residua ordine suo numeri $A',B',C'$ etc. Iam facile intelligitur, esse $x,x^p,x^{pp}$ etc. radices aequationis
\[
z^m-Xz^{m-1}+Yz^{m-2}-Zz^{m-3}+\text{etc.}=0
\]
Quare erit ponendo $z=x$
\[
x^m-Xx^{m-1}+Yx^{m-2}-Zx^{m-3}+\text{etc.}=0
\]
Sed functiones $X-A'$, $Y-B'$, $Z-C'$ etc. per $P$ dividi possunt, quare etiam functio
\[
x^m-A'x^{m-1}+B'x^{m-2}-C'x^{m-3}+\text{etc.}
\]
Hoc autem aliter fieri nequit, nisi sit $A'\equiv A$, $B'\equiv B$, $C'\equiv C$ etc. Q. E. D.

Ceterum notum est, quaecunque alia functio sit $X$ ipsorum $x,x^p$ etc. [in quam omnes hae quantitates aequaliter ingrediantur,] eam semper ex his deduci posse. Ita erit
\SourceLeaf{239}
\[
x^2+x^{2p}+x^{3pp}+\text{etc.}\equiv AA-2B(\mathrm{mod.}\,P)\text{ etc. etc.}
\]
Exempl. Sit $p=5$ et $P\equiv x^2+2x+3$, erit functio $x+x^5$ per $P$ divisa $\equiv-2$, $x^6\equiv3$ etc. etc.
